\setlength{\absparsep}{18pt} 

\begin{resumo}[Resumo]
 \begin{otherlanguage*}{brazil}
   
   \vspace{\onelineskip}
   \noindent 
   
Estudantes da Educação Básica convivem com aplicações de inteligência artificial, mas esse contato cotidiano está longe de assegurar a compreensão de como dados e regras orientam seu funcionamento. Este projeto de qualificação objetiva desenvolver e analisar uma sequência didática apoiada por uma plataforma interativa para o ensino introdutório de \textit{machine learning} no Ensino Médio. O tema entra como conteúdo de Computação, examinando como modelos computacionais partem de dados para gerar previsões numéricas, separar classes ou formar grupos. A proposta articula a Resolução CNE/CEB nº 1/2022 e o pensamento computacional ao uso de múltiplas representações, organizando as aulas da previsão inicial ao experimento prático para sustentar a observação e a explicação. O produto educacional une oito encontros a uma plataforma web executada localmente com dados sintéticos tabulares, cobrindo regressão linear, regressão logística, \textit{k}-vizinhos mais próximos (\textit{k}-NN), árvore de decisão e \textit{k}-means. De natureza aplicada e abordagem qualitativa apoiada por dados descritivos, a pesquisa ocorrerá como intervenção pedagógica em uma turma do próprio pesquisador. O corpus reunirá o diagnóstico inicial e o instrumento final com seu desafio de transferência, somados a três registros intermediários e às notas docentes. Uma rubrica examinará o que muda ou permanece nas dificuldades de relacionar os dados e a tarefa ao efeito de ajustes nas representações, considerando ainda os limites das decisões automatizadas. Frequências e transições entre níveis descreverão o conjunto, sem atribuir à plataforma um efeito causal isolado.
   
   \textbf{Palavras-chave}: Ensino de Computação; Machine Learning; Ensino Médio; Sequência Didática; Plataforma Interativa; Produto Educacional.

 \end{otherlanguage*}
\end{resumo}

% O resumo em inglês (Abstract) será incluído após a validação da versão portuguesa.
